\documentclass[11pt]{article}

\usepackage[margin=1in]{geometry}
\usepackage{amsmath,amssymb,mathtools}
\usepackage{enumitem}
\usepackage[T1]{fontenc}
\usepackage{lmodern}
\usepackage{microtype}

\title{Mathematical Model for D2A Project Draft 2}
\author{}
\date{\today}

\begin{document}
\maketitle

\section{Scope and interpretation}

Draft 2 clears a single selected quarter-hour of submitted supply and demand
bids across the modeled bidding areas. It compares an autarkic case (no
interconnectors) with a capacity-constrained transport network that includes
approximate resistive losses. An optional outer-loop heuristic evaluates new
links and one-step conductor upgrades.

The dispatch model maximizes submitted demand value less submitted supply cost.
It does not include investment costs in this market-clearing objective.
Investment costs are used only to test budget feasibility in the outer search.
Welfare is expressed in EUR/hour because prices are in EUR/MWh and quantities
are in MW. For one 15-minute interval, multiply a welfare value or gain in
EUR/hour by \(0.25\) hours.

This is an illustrative transport model, not an AC or DC optimal power-flow
model. It does not enforce voltage, reactive-power, Kirchhoff-voltage, loop-flow,
unit-commitment, reserve, or reliability constraints.

\section{Sets, inputs, and bid blocks}

Let \(\mathcal{N}\) denote bidding areas and \(\mathcal{E}\) the set of
undirected interconnectors. Each edge \(e\in\mathcal{E}\) has a designated
orientation \(i(e)\to j(e)\) for bookkeeping; the physical flow may use either
direction. Let \(\mathcal{S}_n\) and \(\mathcal{D}_n\) be the supply and demand
blocks for area \(n\).

For each supply block \(s\in\mathcal{S}_n\), let \(q^S_{ns}\) be its maximum
quantity in MW and \(p^S_{ns}\) its offer price in EUR/MWh. For each demand
block \(d\in\mathcal{D}_n\), let \(q^D_{nd}\) and \(p^D_{nd}\) denote its maximum
quantity and bid price, respectively.

The source schedules are cumulative. They are converted into incremental
blocks in source order: the first block quantity is the first cumulative
quantity measured from zero, and each subsequent quantity is the increase over
the previous cumulative quantity. Quantity decreases beyond tolerance are
rejected. Zero-quantity steps are omitted, and adjacent blocks with the same
price (within tolerance) are merged.

\section{Market-clearing variables and objective}

Let \(x^S_{ns}\) be accepted supply and \(x^D_{nd}\) accepted demand, both in
MW. The welfare-maximizing dispatch problem is
\begin{align}
\max_{x^S,x^D,f,\ell,z}\quad
W &=
\sum_{n\in\mathcal{N}}\sum_{d\in\mathcal{D}_n}
    p^D_{nd}x^D_{nd}
-\sum_{n\in\mathcal{N}}\sum_{s\in\mathcal{S}_n}
    p^S_{ns}x^S_{ns}
\label{eq:welfare}\\
\text{subject to}\quad
0\le x^S_{ns} &\le q^S_{ns}
&& \forall n,\ s\in\mathcal{S}_n,
\label{eq:supply-bounds}\\
0\le x^D_{nd} &\le q^D_{nd}
&& \forall n,\ d\in\mathcal{D}_n.
\label{eq:demand-bounds}
\end{align}

Thus accepted demand contributes its bid value and accepted supply incurs its
offer cost. The objective has units EUR/hour and excludes capital expenditure.

\section{Directional flows and nodal balances}

For each edge \(e\), let \(f_e^+\) denote sending flow from \(i(e)\) to
\(j(e)\), and \(f_e^-\) sending flow from \(j(e)\) to \(i(e)\). Let
\(\ell_e^+\) and \(\ell_e^-\) be the corresponding transmission losses, in MW.
The binary variable \(z_e\) selects the permitted direction:
\begin{align}
0\le f_e^+ &\le F_e z_e,\\
0\le f_e^- &\le F_e(1-z_e),\\
z_e &\in \{0,1\}.
\end{align}
Consequently, simultaneous counterflows on a link are prohibited.

For a node \(n\), define its net flow contribution \(g_{ne}\) from edge \(e\)
as
\[
g_{ne} =
\begin{cases}
-f_e^+ + f_e^- - \ell_e^-,
    & n=i(e),\\[2pt]
 f_e^+ - \ell_e^+ - f_e^-,
    & n=j(e),\\[2pt]
0, & n\notin\{i(e),j(e)\}.
\end{cases}
\]
The nodal energy-balance constraints are
\begin{equation}
\sum_{s\in\mathcal{S}_n}x^S_{ns}
-\sum_{d\in\mathcal{D}_n}x^D_{nd}
+\sum_{e\in\mathcal{E}}g_{ne}=0
\qquad \forall n\in\mathcal{N}.
\label{eq:balance}
\end{equation}
An outgoing flow is measured at its sending end; at its receiving end, the
received quantity is sending flow less the associated loss.

\section{Electrical limits and approximate losses}

For each edge, the model uses an estimated length \(L_e\) in km, base conductor
radius \(r_e\), resistivity \(\rho_e\), line-to-line voltage \(V_e\), base
current limit \(I^{\max}_e\), power factor \(\mathrm{PF}_e\), base transfer cap
\(C_e\), and conductor area multiplier \(a_e\). The estimated conductor area,
resistance, current limit, thermal transfer limit, and effective sending-flow
limit are
\begin{align}
A_e &= \pi r_e^2 a_e,\\
R_e &= \frac{\rho_e(1000L_e)}{A_e},\\
I^{\max,\mathrm{up}}_e &= I^{\max}_e a_e,\\
F^{\mathrm{thermal}}_e
&= \frac{\sqrt{3}\,V_e I^{\max,\mathrm{up}}_e\mathrm{PF}_e}{10^6},\\
F_e &= \min\{C_e,F^{\mathrm{thermal}}_e\}.
\label{eq:flow-limit}
\end{align}
Here \(A_e\) is in m\(^2\), \(R_e\) in ohms, \(V_e\) in volts, current in
amperes, and power limits in MW.

For a sending power \(P\) in MW, the balanced three-phase approximation gives
\begin{align}
I_e(P) &= \frac{10^6P}{\sqrt{3}\,V_e\mathrm{PF}_e},\\
\ell_e(P) &= \frac{3I_e(P)^2R_e}{10^6}\\
&= \frac{R_e\,10^6}{V_e^2\mathrm{PF}_e^2}P^2
\qquad \text{MW}.
\label{eq:loss}
\end{align}
The model applies this relation separately to each direction:
\begin{equation}
\ell_e^+ = \widehat{\ell}_e(f_e^+),\qquad
\ell_e^- = \widehat{\ell}_e(f_e^-).
\label{eq:pwl-loss}
\end{equation}
The function \(\widehat{\ell}_e\) is the piecewise-linear interpolant of
\(\ell_e(P)\) at nine evenly spaced points from \(0\) to \(F_e\), giving eight
segments. Inactive-direction flow is zero and has zero modeled loss. This is a
piecewise-linear approximation to the quadratic formula, not an exact AC
network-loss calculation.

\subsection{Illustrative electrical and distance assumptions}

The notebook defaults to aluminium resistivity
\(\rho=2.82\times10^{-8}\ \Omega\mathrm{m}\), conductor radius \(r=0.015\) m,
line voltage \(V=400{,}000\) V, base current limit \(2{,}000\) A, and power
factor \(1\). A conductor upgrade multiplies its area and modeled current limit
by \(2\), holding the base radius fixed.

Market-centre coordinates are used to estimate straight-line distance. For
latitudes \(\phi_i,\phi_j\) and longitudes \(\lambda_i,\lambda_j\), in radians,
the approximation is
\begin{equation}
L_{ij}=6371\sqrt{
(\phi_j-\phi_i)^2+
\left(\cos\frac{\phi_i+\phi_j}{2}\right)^2(\lambda_j-\lambda_i)^2
}\quad\text{km}.
\label{eq:distance}
\end{equation}
This approximation is not routed construction distance. The coordinates,
capacities, and electrical parameters are editable illustrative inputs, not
validated engineering data.

\section{Investment candidates and costs}

The outer search begins from a supplied starting network and evaluates actions
over the one selected market period. Every pair of distinct markets is a
potential build if it is not already linked. A new link has the configured
candidate transfer capacity (by default, the median capacity in the starting
network) and inherits the electrical defaults above.

The new-build cost for pair \((i,j)\), in million GBP, is
\begin{equation}
c^{\mathrm{build}}_{ij}
= \left\lfloor 1.6 L_{ij}\right\rfloor.
\label{eq:build-cost}
\end{equation}
For an existing, not-yet-upgraded link \(e\), the one-step upgrade cost is
\begin{equation}
c^{\mathrm{upgrade}}_e
=0.40(1.6)L_e
=0.64L_e
\quad\text{million GBP}.
\label{eq:upgrade-cost}
\end{equation}
An upgrade sets \(a_e=2\) and doubles the link's base transfer cap. At most one
such upgrade is allowed per link. No action changes the bid curves.

\section{Greedy investment heuristic}

Let \(W(G)\) denote the optimal welfare from
Equations~\ref{eq:welfare}--\ref{eq:pwl-loss} for network \(G\), and let
\(G_k\) be the network after \(k\) accepted actions. The search budget is
\begin{equation}
B = 10\max_{i<j} c^{\mathrm{build}}_{ij}
\quad\text{million GBP},
\label{eq:budget}
\end{equation}
using the configured default budget multiplier of \(10\). With cumulative
spending \(S_k\), the feasible candidate set at iteration \(k\) consists of
unbuilt links and not-yet-upgraded existing links whose cost \(c_a\) satisfies
\[
S_k+c_a\le B.
\]
For every feasible action \(a\), the notebook solves market clearing on the
trial network \(G_k\oplus a\) and computes its one-step incremental welfare:
\begin{equation}
\Delta W_{a,k}=W(G_k\oplus a)-W(G_k).
\end{equation}
It selects the affordable action with the greatest absolute \(\Delta W_{a,k}\),
not the greatest welfare gain per pound:
\begin{equation}
a_k\in\arg\max_{a:\,S_k+c_a\le B}\Delta W_{a,k}.
\end{equation}
The chosen action is accepted only if its gain is greater than the configured
minimum threshold (default \(0.50\) EUR/hour). The search stops when no
affordable candidates remain, the best gain does not exceed the threshold, the
budget is reached, or the iteration limit (default \(50\)) is reached.

This is a one-action-at-a-time greedy heuristic. It does not enumerate all
multi-action portfolios and does not guarantee a globally optimal investment
plan. Capital costs are in GBP while welfare gains are in EUR/hour; the search
uses costs only for budget feasibility and does not perform a discounted
benefit--cost analysis or convert currencies.

\section{Nodal-price diagnostics}

The primary dispatch is a mixed-integer optimization because of the
direction-selection variables. When price diagnostics are requested, the
notebook fixes each edge's selected direction and the piecewise-loss segment
containing its primary-solution flow, then solves a local continuous
piecewise-linear problem. The reported nodal price is derived from the dual
value of the corresponding nodal balance constraint in that local problem.
These are local marginal-price diagnostics conditional on the primary
direction/segment choices, not unrestricted dual prices of the mixed-integer
problem.

\section{Summary of limitations}

\begin{itemize}[leftmargin=*]
\item Only one selected quarter-hour is optimized in a run; outcomes should not
      be interpreted as daily or annual welfare without solving and aggregating
      the corresponding periods.
\item The transport representation does not model full AC power flow, voltage
      variation, reactive power, parallel circuits, or Kirchhoff-voltage laws.
\item Resistivity, coordinates, route lengths, capacities, and investment costs
      are illustrative and should be replaced or validated before engineering or
      empirical conclusions are drawn.
\item The greedy search is path-dependent and myopic. Its selected actions can
      differ from those in a multi-period or globally optimized investment
      formulation.
\item A financial appraisal requires consistent currency, investment timing,
      asset life, discounting, and a benefit horizon.
\end{itemize}

\appendix
\section{Complete algebraic formulation}

This appendix collects the market-clearing problem as one explicit
mixed-integer piecewise-linear program and states how it is embedded in the
multi-start investment experiment. For each market-clearing solve, the network
and its edge parameters are fixed inputs. The outer greedy search changes that
network between solves; investment decisions are not simultaneous decision
variables in the market-clearing program.

\subsection{Index sets and parameters}

\begin{itemize}[leftmargin=*]
\item \(N\): bidding areas; \(E\): undirected links in the network being solved.
\item \(S_n,D_n\): supply and demand blocks at area \(n\).
\item \(q^S_{ns},p^S_{ns}\): maximum supply quantity (MW) and offer price
      (EUR/MWh).
\item \(q^D_{nd},p^D_{nd}\): maximum demand quantity (MW) and bid price
      (EUR/MWh).
\item \(i(e),j(e)\): bookkeeping endpoints for edge \(e\).
\item \(F_e\): effective sending-flow limit (MW), as defined in
      Equation~\ref{eq:flow-limit}.
\item \(K\): number of PWL segments, \(K=8\), with breakpoint index
      \(m=0,\ldots,K\).
\end{itemize}

For each edge, define equally spaced flow breakpoints and their physical loss
values by
\begin{align}
\xi_{em} &= \frac{m}{K}F_e,
&&m=0,\ldots,K,\\
\eta_{em} &= \ell_e(\xi_{em}),
&&m=0,\ldots,K,
\label{eq:breakpoint-values}
\end{align}
where \(\ell_e(\cdot)\) is the quadratic three-phase loss function in
Equation~\ref{eq:loss}. The optimizer uses linear interpolation between
successive points \((\xi_{em},\eta_{em})\).

\subsection{Decision variables}

For every block, \(x^S_{ns}\ge0\) and \(x^D_{nd}\ge0\) are accepted quantities
in MW. For every edge, \(f_e^+,f_e^-\ge0\) and \(\ell_e^+,\ell_e^-\ge0\)
are directional sending flows and losses in MW, and \(z_e\in\{0,1\}\)
selects the direction. For the explicit PWL representation, let
\(\lambda^+_{em},\lambda^-_{em}\ge0\) be interpolation weights for the two
directions.

\subsection{Market-clearing mixed-integer program}

For fixed \(N,E,S_n,D_n\) and edge parameters, the primary dispatch problem is
\begin{align}
\max\quad
&\sum_{n\in N}\sum_{d\in D_n}p^D_{nd}x^D_{nd}
-\sum_{n\in N}\sum_{s\in S_n}p^S_{ns}x^S_{ns}
\label{eq:full-objective}\\
\text{subject to}\quad
&0\le x^S_{ns}\le q^S_{ns}
&&n\in N,\ s\in S_n,
\label{eq:full-supply}\\
&0\le x^D_{nd}\le q^D_{nd}
&&n\in N,\ d\in D_n,
\label{eq:full-demand}\\
&0\le f_e^+\le F_ez_e
&&e\in E,
\label{eq:full-forward-cap}\\
&0\le f_e^-\le F_e(1-z_e)
&&e\in E,
\label{eq:full-reverse-cap}\\
&z_e\in\{0,1\}
&&e\in E,
\label{eq:full-direction}\\
&\sum_{s\in S_n}x^S_{ns}-\sum_{d\in D_n}x^D_{nd}
-\sum_{\substack{e\in E\\n=i(e)}}f_e^+
+\sum_{\substack{e\in E\\n=i(e)}}f_e^-
-\sum_{\substack{e\in E\\n=i(e)}}\ell_e^-
\nonumber\\
&\quad
+\sum_{\substack{e\in E\\n=j(e)}}f_e^+
-\sum_{\substack{e\in E\\n=j(e)}}\ell_e^+
-\sum_{\substack{e\in E\\n=j(e)}}f_e^-=0
&&n\in N,
\label{eq:full-balance}\\
&\sum_{m=0}^{K}\lambda^+_{em}=z_e
&&e\in E,
\label{eq:full-forward-lambda-sum}\\
&\sum_{m=0}^{K}\lambda^-_{em}=1-z_e
&&e\in E,
\label{eq:full-reverse-lambda-sum}\\
&f_e^+=\sum_{m=0}^{K}\xi_{em}\lambda^+_{em},
\qquad
\ell_e^+=\sum_{m=0}^{K}\eta_{em}\lambda^+_{em}
&&e\in E,
\label{eq:full-forward-pwl}\\
&f_e^-=\sum_{m=0}^{K}\xi_{em}\lambda^-_{em},
\qquad
\ell_e^-=\sum_{m=0}^{K}\eta_{em}\lambda^-_{em}
&&e\in E,
\label{eq:full-reverse-pwl}\\
&\lambda^+_{e\cdot}\text{ is SOS2},\qquad
\lambda^-_{e\cdot}\text{ is SOS2}
&&e\in E.
\label{eq:full-sos2}
\end{align}
An SOS2 constraint means that at most two adjacent weights in the ordered
breakpoint list may be nonzero. Equations~\ref{eq:full-forward-lambda-sum}
--\ref{eq:full-sos2} implement the same directional PWL loss graph as the
optimizer's piecewise-linear general constraints: the active direction's
weights interpolate between adjacent breakpoints, while the inactive
direction's weights, flow, and loss are zero. Equations~\ref{eq:full-forward-cap}
--\ref{eq:full-direction} prohibit counterflow. Equation~\ref{eq:full-balance}
withdraws sending flow at its source and injects received flow (sending flow
minus loss) at its destination.

For an empty link set \(E=\varnothing\), the same program reduces to independent
market clearing (autarky), with no edge variables or edge-balance terms.

\subsection{Outer greedy investment procedure}

Let \(G_0\) be a run's starting network and let \(W(G)\) denote the optimal
objective of Equations~\ref{eq:full-objective}--\ref{eq:full-sos2} on network
\(G\). Potential build actions are all absent unordered pairs \(i<j\). A build
uses the median starting-network capacity as its transfer cap and the default
electrical parameters. Its cost is
\[
c^{\mathrm{build}}_{ij}
=\left\lfloor1.6L_{ij}\right\rfloor
\quad\text{million GBP}.
\]
Potential upgrade actions are links whose conductor multiplier remains below
the one allowed upgraded value \(2\). An upgrade sets that multiplier to \(2\),
doubles the link's base transfer cap, and costs
\[
c^{\mathrm{upgrade}}_e=0.64L_e
\quad\text{million GBP}.
\]
The budget for a run is fixed as
\[
B=10\max_{i<j}c^{\mathrm{build}}_{ij}
\quad\text{million GBP}.
\]

At iteration \(k\), let \(G_k\) be the current network and \(C_k\) the set of
available build and upgrade actions. Let \(S_k\) be cumulative spending. The
affordable set is
\[
\mathcal{A}_k=\{a\in C_k:S_k+c_a\le B\}.
\]
For each \(a\in\mathcal{A}_k\), the algorithm forms the trial network
\(G_k\oplus a\), solves the complete market-clearing program, and computes
\[
\Delta W_{a,k}=W(G_k\oplus a)-W(G_k).
\]
If \(\mathcal{A}_k\) is empty, it stops. Otherwise it selects
\[
a_k\in\arg\max_{a\in\mathcal{A}_k}\Delta W_{a,k}.
\]
The action is accepted only when \(\Delta W_{a_k,k}>0.50\) EUR/hour. After
acceptance, \(G_{k+1}=G_k\oplus a_k\) and
\(S_{k+1}=S_k+c_{a_k}\). The search stops when the best increment is at or
below the threshold, the budget is reached, or 50 actions have been accepted.
The selector maximizes absolute welfare gain, not gain per unit of investment
cost; it does not solve a global portfolio optimization.

\subsection{Multi-start design}

The experiment constructs ten connected starting networks over the 13 bidding
areas. From the \(\binom{13}{2}\) possible unordered pairs, it first samples a
spanning tree (12 edges) by applying a seeded random ordering and retaining
edges that connect distinct components, then adds six randomly selected
remaining pairs to obtain 16 links. Existing baseline pairs retain their
configured edge parameters; other pairs receive the default candidate-link
parameters. Candidate seeds are scanned from the fixed seed base, and a
network is accepted only if it shares at most eight edges with every previously
accepted starting network. Thus every pair of accepted starts differs by at
least eight of its 16 edges.

The same selected quarter-hour market data and greedy procedure are solved
independently for each start. For run \(r\), total reported improvement is
\[
\Delta W_r=W(G_{r,\mathrm{final}})-W(G_{r,0}),
\]
while each selected-action gain is marginal relative to that run's current
network at the corresponding iteration. Run comparisons and action-frequency,
price, quantity, and gain-per-cost summaries are descriptive; none are optimized
across runs. Ten different seeds explore starting-topology sensitivity but do
not establish a global optimum or results for other delivery periods.

\end{document}
